% ============================================================
% Section 53: N 过程与 U 过程——声子热平衡与热阻的机制
% ============================================================
\newpage
\section{固体理论：N 过程与 U 过程——声子热平衡与热阻的机制}
\label{sec:nu-process-thermal-equilibrium}

\begin{modelbox}[题目：三声子碰撞与热平衡的建立]
非金属结晶固体的热传导需要一种使声子分布趋于热平衡的机制。
\begin{enumerate}
    \item 说明 $\vec q_1+\vec q_2=\vec q_3$ 型的三声子碰撞过程不能建立热平衡；
    \item 描述建立热平衡过程的形式。
\end{enumerate}
\end{modelbox}

\begin{tipbox}[考点快速分析]
\begin{itemize}
    \item \textbf{N 过程}（Normal Process）：$\vec q_1+\vec q_2=\vec q_3$，总准动量守恒，$\vec G=0$
    \item \textbf{U 过程}（Umklapp Process）：$\vec q_1+\vec q_2=\vec q_3+\vec G$，总准动量被倒格矢改变，$\vec G\neq0$
    \item \textbf{热流与准动量}：热流密度 $\vec J$ 与声子总准动量 $\sum\hbar\vec q$ 方向一致
    \item \textbf{热阻来源}：只有能改变总准动量的过程才能产生热阻，使系统弛豫到热平衡
\end{itemize}
\end{tipbox}

\begin{derivationbox}[标准解题过程]
\textbf{Step 0: 三声子过程的准动量守恒}

在非简谐相互作用下，晶体中可发生三声子碰撞：两个声子合并为一个，或一个声子分裂为两个。准动量守恒的一般形式为
\begin{equation}
\vec q_1+\vec q_2=\vec q_3+\vec G,
\end{equation}
其中 $\vec G$ 为任一倒格矢（包括零矢量）。根据 $\vec G$ 是否为零，分为两类：

\begin{itemize}
    \item \textbf{N 过程}（$\vec G=0$）：$\vec q_1+\vec q_2=\vec q_3$，总准动量严格守恒
    \item \textbf{U 过程}（$\vec G\neq0$）：$\vec q_1+\vec q_2=\vec q_3+\vec G$，总准动量被倒格矢改变
\end{itemize}

\textbf{Step 1: N 过程不能建立热平衡}

声子系统的总准动量为
\begin{equation}
\vec P_{\text{total}}=\sum_{\vec q,s}\hbar\vec q\,n_s(\vec q).
\end{equation}

热流密度与声子分布的关系：
\begin{equation}
\vec J=\sum_{\vec q,s}\hbar\omega_s(\vec q)\,\vec v_s(\vec q)\,n_s(\vec q),
\label{eq:heat-current}
\end{equation}
其中群速度 $\vec v_s(\vec q)=\nabla_{\vec q}\omega_s(\vec q)$。在德拜模型中 $\vec v_s\parallel\vec q$，故热流方向与总准动量方向一致。

N 过程保持总准动量不变：
\begin{equation}
\vec P_{\text{total}}^{\text{初}}=\vec P_{\text{total}}^{\text{末}}.
\end{equation}

这意味着 N 过程只能使声子在保持总准动量的前提下重新分布——声子之间交换能量和个体动量，但不改变整个声子气体的``漂移速度''。若初始存在定向热流（$\vec P_{\text{total}}\neq0$），N 过程无法将其消除，热流将持续存在，热导率趋于无穷大。

\textbf{结论}：仅有 N 过程不能使声子分布弛豫到总准动量为零的热平衡态。

\textbf{Step 2: 建立热平衡的机制——U 过程}

U 过程中倒格矢 $\vec G\neq0$ 介入，碰撞前后的总准动量改变：
\begin{equation}
\Delta\vec P_{\text{total}}=\hbar\vec G\neq0.
\end{equation}

这部分准动量被转移给了晶格整体（宏观上表现为边界或支撑结构吸收）。U 过程的物理作用是：

\begin{enumerate}
    \item \textbf{破坏定向热流}：U 过程可显著改变甚至反转声子波矢方向，有效耗散总准动量
    \item \textbf{提供热阻}：使系统在有限弛豫时间内达到总准动量为零的热平衡分布，产生有限热导率
    \item \textbf{高温主导}：U 过程需要参与声子的波矢足够大（$\vec q_1+\vec q_2$ 超出布里渊区边界），因此需要高频声子。高温下（$T\gtrsim\Theta_D$）大量高频声子被激发，U 过程频繁，成为热阻的主要来源
\end{enumerate}

\textbf{低温下的补充}：极低温（$T\ll\Theta_D$）下高频声子极少，U 过程极难发生。此时热阻主要由晶体边界、杂质、同位素等缺陷散射提供，这些过程同样不守恒准动量。
\end{derivationbox}

\begin{formulabox}[答案汇总]
\begin{center}
\begin{tabular}{p{3.5cm}p{8cm}}
\toprule
问题 & 答案 \\
\midrule
(1) N 过程为何不能建立热平衡 & N 过程保持声子总准动量守恒，无法耗散定向热流，不能产生热阻，热导率趋于无穷大 \\
\addlinespace
(2) 建立热平衡的机制 & U 过程：$\vec q_1+\vec q_2=\vec q_3+\vec G$（$\vec G\neq0$）。通过倒格矢改变总准动量，耗散定向热流，使声子分布弛豫到热平衡，产生有限热阻 \\
\bottomrule
\end{tabular}
\end{center}
\end{formulabox}

\begin{insightbox}[物理图像与概念深化]
\textbf{1. N 过程 vs U 过程的本质区别——以气体类比}

\begin{center}
\begin{tabular}{p{2.5cm}p{5.5cm}p{5.5cm}}
\toprule
 & \textbf{N 过程} & \textbf{U 过程} \\
\midrule
碰撞规则 & 分子间弹性碰撞，总动量守恒 & 分子与器壁碰撞，动量被壁吸收 \\
\addlinespace
气流影响 & 分子交换动量但气流方向不变 & 分子反向弹回，气流被破坏 \\
\addlinespace
热阻类比 & 理想气体无黏性（无热阻） & 黏性/摩擦导致热阻 \\
\bottomrule
\end{tabular}
\end{center}

\textit{图像}：把声子气体想象成管道里的气体分子。N 过程是分子之间的碰撞——总动量不变，气体整体仍沿管道流动。U 过程是分子与粗糙管壁的碰撞——分子被弹回甚至反向，整体流动被破坏，形成阻力。

\textbf{2. 为什么 U 过程需要大波矢声子？}

U 过程要求 $\vec q_1+\vec q_2$ 超出第一布里渊区边界，即 $|\vec q_1+\vec q_2|>\pi/a$。这意味着至少有一个参与声子的波矢要接近布里渊区边界（$q\sim\pi/a$），对应频率 $\omega\sim\omega_D$。而普朗克分布中，频率为 $\omega_D$ 的声子占据数为
\begin{equation}
n(\omega_D)=\frac{1}{e^{\Theta_D/T}-1}.
\end{equation}

\begin{itemize}
    \item 高温（$T\gtrsim\Theta_D$）：$n(\omega_D)\sim O(1)$，大量高频声子存在，U 过程频繁
    \item 低温（$T\ll\Theta_D$）：$n(\omega_D)\approx e^{-\Theta_D/T}\to0$，高频声子指数稀少，U 过程被``冻结''
\end{itemize}

这就是 U 过程具有强烈温度依赖性的根本原因。

\textbf{3. 热导率的温度依赖——两段论}

非金属晶体的热导率 $\kappa$ 随温度的变化呈现典型的``倒 U 型''：

\begin{itemize}
    \item \textbf{高温区}（$T\gtrsim\Theta_D$）：U 过程主导，声子平均自由程 $l\propto1/T$，故 $\kappa\propto1/T$
    \item \textbf{中温区}（$T\sim\Theta_D/10$）：U 过程减少但仍有贡献，$l$ 增大，$\kappa$ 出现峰值
    \item \textbf{低温区}（$T\ll\Theta_D$）：U 过程冻结，边界/缺陷散射主导，$l\sim$ 晶体尺寸（Casimir 极限），$\kappa\propto T^3$
\end{itemize}

峰值温度附近的热导率是材料纯度和晶体质量的敏感指标——高纯度单晶在低温下可达极高的热导率（如蓝宝石在 $30\,\mathrm{K}$ 时热导率超过铜）。

\textbf{4. 与 sec50 的关联}

sec50（三声子过程的具体计算）和 sec53（N/U 过程的物理后果）是一个整体：
\begin{itemize}
    \item sec50：展示 U 过程的``几何''——准动量守恒 + 布里渊区折叠，计算末态波矢
    \item sec53：展示 U 过程的``物理''——为什么必须 U 过程才能产生热阻，N 过程为何不足
\end{itemize}

从计算到物理的完整链条：先学会判断一个给定的三声子碰撞是 N 还是 U，再理解为什么只有 U 过程能使系统达到热平衡。
\end{insightbox}

\begin{thinkbox}[考场速记：N/U 过程与热平衡问题的答题框架]
遇到``为什么 N 过程不能建立热平衡''或``热阻的来源是什么''类问题：
\begin{enumerate}
    \item \textbf{写准动量守恒}：$\vec q_1+\vec q_2=\vec q_3+\vec G$，区分 N（$\vec G=0$）和 U（$\vec G\neq0$）
    \item \textbf{说明 N 过程的问题}：总准动量守恒 → 无法耗散定向漂移 → 热导率发散
    \item \textbf{说明 U 过程的作用}：倒格矢吸收准动量 → 破坏定向热流 → 提供热阻
    \item \textbf{补充温度依赖性}：高温 U 过程活跃，低温 U 过程冻结（需要高频声子）
\end{enumerate}
\textit{核心逻辑链}：热平衡 $\Leftrightarrow$ 总准动量为零 $\Leftrightarrow$ 需要改变总准动量的机制 $\Leftrightarrow$ 只有 U 过程能做到。
\end{thinkbox}
